%% LyX 1.3 created this file.  For more info, see http://www.lyx.org/.
%% Do not edit unless you really know what you are doing.
\documentclass[11pt,a4paper,english]{article}
\usepackage[latin1]{inputenc}
\setcounter{secnumdepth}{3}
\setcounter{tocdepth}{3}
\IfFileExists{url.sty}{\usepackage{url}}
                      {\newcommand{\url}{\texttt}}

\makeatletter

%%%%%%%%%%%%%%%%%%%%%%%%%%%%%% LyX specific LaTeX commands.
%% Bold symbol macro for standard LaTeX users
\newcommand{\boldsymbol}[1]{\mbox{\boldmath $#1$}}

%% Because html converters don't know tabularnewline
\providecommand{\tabularnewline}{\\}

%%%%%%%%%%%%%%%%%%%%%%%%%%%%%% Textclass specific LaTeX commands.
 \usepackage{pslatex}
 \usepackage{myexam}
 \usepackage{graphicx}
 \usepackage{amsmath}
 \usepackage{epic}
 \usepackage{eepic}
 \input{neilutils}
 \newenvironment{lyxcode}
   {\begin{list}{}{
     \setlength{\rightmargin}{\leftmargin}
     \setlength{\listparindent}{0pt}% needed for AMS classes
     \raggedright
     \setlength{\itemsep}{0pt}
     \setlength{\parsep}{0pt}
     \normalfont\ttfamily}%
    \item[]}
   {\end{list}}

%%%%%%%%%%%%%%%%%%%%%%%%%%%%%% User specified LaTeX commands.
\newcommand{\answer}[2]{}
\newcommand{\masters}[2]{#2}

\usepackage{babel}
\makeatother
\begin{document}

\pressmark{\masters{COM6131}{COM3270}}


\examtitle{NETWORK PERFORMANCE ANALYSIS}


\examdate{AUTUMN SEMESTER 2002--2003}


\examtime{2 HOURS}

\begin{exam}
~
\begin{question}
In an $M/G/1$ queue inter-arrival times are exponentially distributed
and service times are `generally' distributed. 
\begin{qupart}
Use the concept of residual service time to show that, assuming customers
are served in the order that they arrive, the average waiting time
is given by\[
\bar{W}=\frac{\bar{R}}{1-\rho},\]
where $\bar{R}$ is the average residual service time and $\rho$
is the queue utilization factor. Use diagrams, where necessary, to
explain your reasoning and state your assumptions clearly.

\mypercent{50}

\end{qupart}
~

\begin{qupart}
Use a graphical argument to derive a formula for the residual service
time, then combine this equation with the formula for the average
waiting time above to show that\[
\bar{W}=\frac{\lambda\left\langle X^{2}\right\rangle }{2\left(1-\rho\right)},\]
where $\lambda$ is the arrival rate and $\left\langle X^{2}\right\rangle $
is the expectation of the squared service time.

\mypercent{40}

\end{qupart}
~

\begin{qupart}
The equation you have derived is easily recognized as the Pollaczek-Khinchin
formula. Use it to derive the average waiting time in an $M/D/1$
queue.

\mypercent{10}

\end{qupart}
\end{question}
~\turnover

\begin{question}
You are presented with the following portion of Tcl code designed
for the network simulator NS2.
\begin{lyxcode}
set~node1~{[}\$ns~node{]}~\\
set~node2~{[}\$ns~node{]}~\\
set~node3~{[}\$ns~node{]}~\\
set~node4~{[}\$ns~node{]}~\\
~\\
\$node1~set~multiPath\_~1~\\
~\\
set~link1~{[}\$ns~duplex-link~\$node1~\$node2~1Kb~0ms~DropTail{]}~\\
set~link2~{[}\$ns~duplex-link~\$node1~\$node3~1Kb~0ms~DropTail{]}~\\
set~link3~{[}\$ns~duplex-link~\$node2~\$node4~1Kb~0ms~DropTail{]}~\\
set~link4~{[}\$ns~duplex-link~\$node3~\$node4~1Kb~0ms~DropTail{]}~\\
~\\
set~src~{[}new~Agent/UDP{]}~\\
\$ns~attach-agent~\$node1~\$src~\\
set~sink~{[}new~Agent/Null{]}~\\
\$ns~attach-agent~\$node4~\$sink~\\

\end{lyxcode}
\begin{qupart}
Assuming that only the specified nodes exist, draw the network topology
showing the associated routes.

\mypercent{25}

\end{qupart}
~

\begin{qupart}
State the assumptions behind the Kleinrock Independence approximation.
Discuss to what extent they are valid for this network topology.

\mypercent{35}

\end{qupart}
~

\begin{qupart}
Packet arrivals occur as a Poisson distribution at \url{node1} with
a rate of 2 per second. Packet lengths are distributed as an exponential
distribution with a mean of 0.5 Kb. Assuming that packets are sent
randomly along all possible routes with equal probability, use Kleinrock's
independence assumption to estimate the steady state average number
of packets that are at each link.

\mypercent{25}

\end{qupart}
~

\begin{qupart}
What is the expected delay for each route through the network?
\end{qupart}

\mypercent{15}

\end{question}
~\continued

\begin{question}
\textbf{Hand in the Erlang C curves provided with your answer to this
question.}

Northern Telecom has just completed the purchase of several local
telecommunications companies, each with their own regional support
teams. There are three regional support teams. The number of staff
in each team and the call rate for each team's centre is given in
the following table. You can assume that calls arrive as a Poisson
process and that the service times are exponentially distributed.

\begin{center}\begin{tabular}{|c|c|c|c|}
\hline 
Region&
A&
B&
C\tabularnewline
\hline
\hline 
$m$&
10&
13&
10\tabularnewline
\hline
\hline 
customers/hour&
120&
160&
100\tabularnewline
\hline
\end{tabular}\end{center}

\begin{qupart}
Given that the average holding time, $\tau$, for support queries
is 3 minutes, use the Erlang C curves at the back of the paper to
give, for each region, the probability of queueing, $C\left(m,a\right)$. 

\mypercent{30}

\end{qupart}
~

\begin{qupart}
The formula for the average queue length in an $M/M/m$ system is\[
\bar{N}_{Q}=\frac{\rho}{1-\rho}C\left(m,a\right)\]
Use Little's formula to compute the formula for the average queueing
time, and hence compute the expected queueing time for each call centre.

\mypercent{30}

\end{qupart}
~

\begin{qupart}
Given that the average queueing time for customers who \emph{don't}
queue is zero. What is the average queueing time for customers who
\emph{do} queue?

\mypercent{20}

\end{qupart}
~

\begin{qupart}
The company wishes to cut costs by merging the call centres to one
location. Assuming that the total traffic level, the average holding
time and the total staff numbers stay the same what will be the average
queueing time for the new centre? 

\mypercent{20}

\end{qupart}
\end{question}
~

\turnover

\begin{question}
A call centre employs thirty operators selling tickets for a major
sporting event. The call centre has sixty lines in total. When all
operators are active new calls are placed in a queue for the first
available operator. The call centre is to be modelled as an $M/M/m/n$
queue, \emph{i.e.} a queue in which there are $m$ servers and a maximum
state of $n$.
\begin{qupart}
Draw a state diagram representing the queueing system.

\mypercent{30}

\end{qupart}
~

\begin{qupart}
Write down the detailed balance equations for the two cases where
the queue state $k<m$ and when $m\leq k<n$.

\mypercent{20}

\end{qupart}
~

\begin{qupart}
For each of the two cases outlined in the previous part write down
the steady state distribution as a function of $P\left(0\right)$,
the probability of the system being empty. 

\mypercent{30}

\end{qupart}
~

\begin{qupart}
Use the fact that the steady state distribution must be normalized
to show that \[
P\left(0\right)=\left[\sum_{k=0}^{m}\frac{a^{k}}{k!}+\sum_{k=m+1}^{n}\frac{a^{k}}{m!m^{k-m}}\right]^{-1}\]
where $a=\frac{\lambda}{\mu}$, $\lambda$ is the arrival rate and
$\mu$ is the rate of the service time distribution.

\mypercent{20}

\end{qupart}
\end{question}
~

\questionsend

\myfooter{}\newpage\myheader

\topmargin -1.8cm \raggedleft{{\bf Registration Number \rule{4cm}{.5pt}}}\newcommand{\figureTextSize}{\tiny}

\begin{center}\input{diagrams/erlangCCurves0to20.eepic}\end{center}\end{exam}

\end{document}
